\chapter{Homotopy}

\noindent\textbf{General Idea:} To understand a space $X$, we can study continuous maps $f: A \rightarrow X$ or $g: X\rightarrow B$. For example, $A = [0,1] \leftrightarrow$ path-connected; $A = \mathbb{S}^1 \leftrightarrow$ loops.

\begin{definition}{Homotopy}\\
Let $X,Y$ be two topological spaces. A homotopy between two continuous maps $f,g :X \rightarrow Y$ is a continuous map $H : X\times [0,1] \rightarrow Y$, s.t.:
$$H(x,0) = f(x) \;\;\;\;and\;\;\;\;  H(x,1) = g(x).$$
If such $H$ exists, we say $f$ and $g$ are homotopic, denoted as
$$f \simeq g.$$
\end{definition}

\noindent\textbf{Example:} $Y= \mathbb{R}^2$. Consider two arbitrary cont maps $f,g: X \rightarrow Y$. Then define:
$$H(x,t) := (1-t) f(x) + tg(x),$$
an interpolate. The well-defined-ness of $H$ is by the convexivity of $\mathbb{R}^2$, i.e.: for all $x,t$, $H(x,t) \in Y$.

\begin{proposition}
    Homotopy is an equivalence relation.
\end{proposition}

\begin{proof}
    \begin{itemize}
        \item Take $H(x,t) = f(x)$ for all $t$. Then $f \simeq f$.
        \item If $f \simeq g$, then there exists $H$ in the definition of homotopy. Then define $H': X \times [0,1] \rightarrow Y$ by $H' (x,t) := H(x,1-t)$. Hence, $g \simeq f$.
        \item It $f\simeq g$ and $g \simeq f$, then there are two continuous maps $H:X \times [0,1] \rightarrow Y$ and $K: X\times [0,1] \rightarrow Y$ with $H(x,0) = f(x)$, $H(x,1) = K(x,0) = g(x)$, and $K(x,1)=h(x)$. Then define
        $$J(x,t) := \begin{cases}
            H(x,2t), &0\leq t \leq \frac{1}{2}\\
            K(x,2t-1) &\frac{1}{2} \leq t \leq 1
        \end{cases}$$
        It's clear that $J$ is cont. Then $f \simeq h$.
    \end{itemize}
\end{proof}

\begin{definition}{Homotopy Class}\\
Let $f: X \rightarrow Y$ be a continuous map. A homotopy class $[f]$ is the collection of continuous maps homotopic to $f$:
$$[f] := \{ g: X \rightarrow Y \mid g \simeq f \}.$$
\end{definition}

\noindent\textbf{Example:} If $Y = \mathbb{R}^2$, then $[f] = \mathcal{C}(X \rightarrow \mathbb{R}^2)$.

\begin{definition}{Homotopic Equivalence}
Anyway a tempt for learning github.
\end{definition}